\documentclass[a4paper,12pt]{article} 
\usepackage[utf8]{inputenc} 
\usepackage[T1]{fontenc} 
\usepackage[french]{babel} % Police Times 
\usepackage{newtxtext,newtxmath}
\usepackage[a4paper,margin=2.5cm]{geometry} 
\usepackage{fancyhdr} 
\pagestyle{fancy} 
\fancyhead[L]{Philippe-Cadet GANTIN} 
\fancyhead[R]{Génie Logiciel} 
\fancyfoot[C]{\thepage} 
\begin{document} 
	\section{Exercice} 
	L'UNSTIM a enregistré quatre cas de vol de motos. Après réflexion, le président des étudiants de l'université a lancé un hackathon consistant à mettre en place un projet de sécurisation des motos via un démarrage biométrique. Après avoir procédé à une analyse complète du dispositif, réaliser la modélisation des diagrammes suivants : 
	\begin{itemize} 
		\item Diagramme des cas d'utilisation ; 
		\item Diagramme de séquence ; 
		\item Diagramme d'activité ; 
		\item Diagramme des composants. 
	\end{itemize} 
	
	\section{Résultat} 
	
	\subsection{Analyse complète du dispositif} 
	\textbf{Mots clés}
	
	\textbf{Proposition de thème}
	
	\textbf{Composants du dispositif}
	
	\textbf{Schéma de fonctionnement de dispositif}
	
	\textbf{Scénario pratique possible}
	\subsection{Diagramme des cas d'utilisation}
	\subsection{Diagramme de séquence}
	\subsection{Diagramme des composants}
\end{document}